\chapter{Primitive-Recursive Texts and Tester Identification}
\label{sec:p15-pr-learning-obstructions}

Restricting a positive presentation to a primitive-recursive sequence
does not by itself make exact tester identification possible.
Gold's Theorem I.9 separates this restriction from the generator-naming
result in \cref{sec:p15-formal-identification}
\cite[pp.\ 471--473]{Gold1967IdentificationFLT}.
The obstruction needs a primitive-recursive guessing rule, not merely
a computable rule. Gold invokes Theorem~5 of his 1965 limiting-recursion
paper for that reduction \cite{Gold1965LimitingRecursionPR}.
The finite-prefix normalization below is supplied as an explicit
serial simulation, with its own precisely declared scope.

\section{The typed interfaces and finite computations}
\label{sec:p15-pr-learning-interfaces}

\begin{hbconvention}{def:p15-pr-learning-codes}{Word, text, tester and computation codes}
Use the fixed primitive-recursive word coding of
\cref{def:p15-flt-names}, shifted if necessary so words have positive
integer codes. In this chapter a text is
\(i_1,i_2,\ldots\), and a finite observation prefix is
\(\sigma=(i_1,\ldots,i_t)\).
The text function \(t\mapsto i_t\), for \(t\geq1\), must be
primitive recursive. Repetitions are allowed and its range is the
target language. There are no negative labels or pause symbols.

Fix one effective machine-program numbering throughout the argument.
Gold's \(Z_a(x)\) denotes the partial numerical output of the program
with hypothesis number \(a\), on word code \(x\).
A valid \emph{tester name} for \(L\) must be total on word codes, have values in
\(\{0,1\}\), and satisfy \(Z_a(x)=\mathbf1_L(x)\) on every word.
For this tester-naming instance of the earlier \(\Gamma\),
\(\Gamma(a)=L\) exactly when \(a\) is such a valid tester for \(L\);
on other program numbers the naming map is undefined.
An intermediate hypothesis may be a wrong tester or a partial program.
No uniform procedure decides which program numbers are total testers.
A tester-identifiable target class must be contained in this
tester-naming range, as required by \cref{def:p15-flt-names}.
An unnameable target is a failure of that prerequisite, not success
inferred by vacuous quantification over absent presentations.
A learner is a total computable prefix function
\(G_0(\sigma)\in\mathbb N\); an optional abstention can instead be
encoded by a fixed invalid program number.
Success requires eventual constancy of that actual number and
correctness of its tester meaning.

Keep Gold's native computation notation
\[
 T(a,x,y),\qquad U(y).
\]
Here \(T\) is a numerical three-place primitive-recursive predicate:
\(y\) codes a complete halting computation of program \(a\) on input
\(x\). The primitive-recursive numerical function \(U\) extracts the
output of such a computation. These are not a presentation sequence
called \(T\), a semantic truth predicate, or an oracle for termination.
The coding is chosen consistently with \(Z_a\); malformed histories
make \(T\) false, and \(U\) has an arbitrary total default on them.
For valid histories,
\[
 (\exists y)\,[T(a,x,y)\land U(y)=v]
 \quad\Longleftrightarrow\quad Z_a(x)\downarrow=v.
 \tag{PR.1}
\]
Thus a bound \(y\leq t\) limits the \emph{history code}, not necessarily
the number of computation steps.
The symbols retain the roles of Gold's equation I.16.
In (PR.3), Gold's \(\equiv\) defines a predicate and his
\(\&\) is Boolean conjunction. This local definitional use of
\(\equiv\) is not elementary equivalence between structures.
\end{hbconvention}

\begin{hblemma}{thm:p15-pr-learning-coding}{proved=annex:p15-pr-learning-coding}{Primitive-recursive finite simulation}
There is a standard finite-sequence and machine-configuration coding
with primitive-recursive prefix extraction, initialization,
one-step transition, halting recognition and numerical output
extraction. Complete finite-history verification gives the \(T,U\)
of \cref{def:p15-pr-learning-codes} and satisfies (PR.1).
Iterating a fixed transition for an explicitly given finite number
of steps is primitive recursive. This is not a primitive-recursive
universal evaluator that runs every terminating program to completion.
\end{hblemma}

\section{Normalization without losing a hypothesis change}
\label{sec:p15-pr-learning-normalization}

\begin{hbmethod}{def:p15-pr-learning-relay}{The serial prefix simulator}
Fix a machine computing \(G_0\) on coded finite sequences.
For an input \(\sigma=(i_1,\ldots,i_n)\), start with next-prefix
number \(m=1\), an inactive simulation state, and a fixed initial
hypothesis \(a_{\mathrm{init}}\).
Perform exactly \(n\) ticks, labelled \(1,\ldots,n\).
At a tick, initialize the inactive machine on
\(\sigma{\upharpoonright}m\), if necessary, and execute one
machine transition. If it is now halted, copy its numerical
answer into the retained hypothesis, increase \(m\) by one and
mark the simulation inactive for the next tick.
If it has not halted, retain its configuration and the previous
hypothesis. At most one prefix computation is completed per tick,
including an initially halted computation.
Return the retained hypothesis; call this prefix function
\(\widehat G(\sigma)\). On the empty prefix return
\(a_{\mathrm{init}}\).
On a numerical input that is not the code of a finite sequence
of word numbers, return the same default.

Before tick \(r\), at most \(r-1\) computations have completed, so
\(m\leq r\leq n\). Every prefix requested exists in the input.
Initialization, a transition, copying, and resetting are finite
primitive-recursive operations. No unbounded waiting occurs inside
a tick. A long-running call is spread across successive ticks.
\end{hbmethod}

\begin{hbtheorem}{thm:p15-pr-learning-normalization}{proved=annex:p15-pr-learning-normalization}{Exact finite-prefix convergence normalization}
For every total computable finite-prefix rule \(G_0\), the
\(\widehat G\) of \cref{def:p15-pr-learning-relay} is primitive
recursive on finite-sequence codes. For every infinite sequence
\(i=(i_1,i_2,\ldots)\) and every hypothesis number \(a\),
\[
 \bigl[G_0(i_1,\ldots,i_t)\longrightarrow a\bigr]
 \quad\Longleftrightarrow\quad
 \bigl[\widehat G(i_1,\ldots,i_t)\longrightarrow a\bigr].
 \tag{PR.2}
\]
Arrows here mean eventual equality of natural numbers, not a metric
limit. In particular nonstabilization is preserved, and a stable
tester number is not replaced by a different equivalent name.
The construction does not change the observed presentation,
and applies to primitive-recursive texts as well as to arbitrary
ones. It asserts no uniform stabilization time or runtime bound.
\end{hbtheorem}

\begin{hbexample}{def:p15-pr-learning-relay-trace}{A slow computation cannot be skipped}
Suppose, on a particular text, the first five ordinary prefix
answers are \(7,7,11,7,7\), and their simulations take respectively
\(2,4,1,3,2\) ticks, under the fixed transition convention.
The serial simulator's answers at ticks \(1\) through \(12\) are
\[
 a_{\mathrm{init}},7,7,7,7,7,11,11,11,7,7,7.
\]
The answer \(11\) is displayed when its call finishes at tick \(7\),
even though later input is already available.
If answers \(7\) and \(11\) alternate forever, both remain visible
infinitely often. Selecting only the largest prefix whose call
finishes within the current budget can instead skip slow calls;
that shortcut does not establish the reverse implication in (PR.2).
The displayed numbers are illustrative hypothesis codes, not
membership values of the text's words.
\end{hbexample}

\section{The primitive-recursive countertext}
\label{sec:p15-pr-learning-countertext}

\begin{hblemma}{thm:p15-pr-learning-range}{proved=annex:p15-pr-learning-range}{A primitive-recursive range for a nonempty enumerable language}
Every nonempty computably enumerable language \(L\), in the declared
word coding, is the range of a total primitive-recursive function
\(f:\mathbb N_{>0}\to\mathbb N_{>0}\).
The construction may repeat words and uses one fixed word already
known to belong to \(L\). It is not a uniform decision of emptiness
and does not assert that membership in \(L\) is primitive recursive.
\end{hblemma}

\begin{hbmethod}{def:p15-pr-learning-native-diagonal}{Gold's bounded rejection construction}
Let \(\mathcal C\) contain every finite language and an infinite
language \(L\). If \(L\) is not recursive it has no tester name,
so tester identification of \(\mathcal C\) already fails.
Otherwise fix a primitive-recursive \(f\) with range \(L\).
For a proposed successful learner \(G_0\), use
\(\widehat G\) from (PR.2), and denote this primitive-recursive
prefix rule by \(G\), retaining Gold's native notation.
Define simultaneously
\[
 \begin{split}
 g_t&=G(i_1,\ldots,i_t),\\
 P_t&\equiv[f(X_t)=i_1]\lor\cdots\lor[f(X_t)=i_t],\\
 Q_t&\equiv(\exists y\leq t)
       \{T[g_t,f(X_t),y]\mathbin{\&}[U(y)=0]\},\\
 X_1&=1,\qquad i_1=f(1),\\
 X_{t+1}&=
 \begin{cases}X_t+1,&P_t,\\X_t,&\neg P_t,\end{cases}\\
 i_{t+1}&=
 \begin{cases}f(X_t),&\neg P_t\land Q_t,\\i_t,&\text{otherwise}.
 \end{cases}
 \end{split}
 \tag{PR.3}
\]
These are Gold's equations I.14--I.22
\cite[pp.\ 471--472]{Gold1967IdentificationFLT}.
The information prefix, current hypothesis and candidate word all
have different types. In \(T[g_t,f(X_t),y]\), the first argument
is a tester program number, the second a word number and the
third a computation-history number.

\(P_t\) says the current candidate is already in the observed
positive content. If it is, advance its enumeration position.
Otherwise keep the same position until bounded evidence appears
that the current hypothesis rejects that candidate; then present it.
When evidence is absent, repeat the last positive word.
The predicates are bounded, and the paired sequence is defined
by primitive recursion on its finite coded history.
\end{hbmethod}

\begin{hbtheorem}{thm:p15-pr-learning-obstruction}{proved=annex:p15-pr-learning-obstruction}{Gold I.9: superfinite classes resist primitive-recursive-text tester learning}
For information given by primitive-recursive positive texts and
the tester-naming relation, no class containing all finite languages
and at least one infinite language is identifiable in the limit
by a total computable finite-prefix learner.
When the selected infinite member \(L\) is recursive,
for each proposed learner (PR.3) produces a primitive-recursive
countertext. Its range is either a finite nonempty subset of \(L\)
or all of \(L\). If the learner's hypotheses stabilize, the
stabilized program fails the tester specification for that range;
if they do not stabilize, identification already fails.
If \(L\) is nonrecursive, the tester-name prerequisite itself fails,
as specified in \cref{def:p15-pr-learning-codes}.
\end{hbtheorem}

\begin{hbexample}{def:p15-pr-learning-candidate-case}{What the delayed word proves}
Display word names in this example; each denotes its fixed positive
integer code. Here \(a\) is an alphabet letter, not a hypothesis number.
For the unary-word language \(\{\varepsilon,a,a^2,\ldots\}\),
take \(f(k)=a^{k-1}\). Initially \(i_1=\varepsilon\), \(X_1=1\);
because the first candidate was seen, \(X_2=2\) and
\(i_2=\varepsilon\). The next candidate \(a\) is now held while
\(\varepsilon\) is repeated. If a current tester is observed,
through a finite computation witness, to reject \(a\), then \(a\)
is inserted and that hypothesis is wrong for every eventual
target containing \(a\). No fixed number of repetitions forces
a rejection; only the witnessed computation triggers insertion.
The same mechanism next applies to \(a^2\), with any delays
determined by the learner and the computation-history numbering.
\end{hbexample}

\begin{proofboundary}{def:p15-pr-learning-limits}{Which assumptions do the work}
The negative result is about tester output and eventual syntactic
stabilization, not merely eventual semantic correctness of changing
names. Under generator naming, primitive-recursive texts have a
different positive result, as in \cref{thm:p15-flt-pr-text}.
An accepting-only recognizer cannot certify the rejection required
by \(Q_t\). Totality is required of a correct final tester, not of
all intermediate guesses, and no totality oracle is introduced.
Each constructed presentation is primitive recursive for the fixed
proposed learner; this is not a uniform primitive-recursive
evaluation of every computable learner.

Nonempty pure texts suffice for the contradiction, so adding an
empty-language pause convention does not repair it.
Negative data, a restricted target family omitting some finite
sublanguages, or a different hypothesis-success criterion changes
the theorem's premises. The finite-prefix normalization is distinct
from a general theorem about numerical/function-variable limiting
functionals; neither provides a finite stopping certificate.
\end{proofboundary}
